\documentclass[12pt,letterpaper]{article}

% Pacotes básicos
\usepackage[utf8]{inputenc}
\usepackage[T1]{fontenc}
\usepackage[brazil]{babel}
\usepackage{amsmath, amssymb, amsthm}

% Estilo de teorema
\newtheorem{theorem}{Teorema}

% Espaçamento
\usepackage{setspace}
\onehalfspacing

% Margens
%\usepackage[margin=1in]{geometry}
\usepackage{fullpage}
\usepackage{enumitem}

\begin{document}

\begin{center}
  {\LARGE Propriedade Forte de Markov}
\end{center}

\bigskip

\begin{theorem}[Propriedade Forte de Markov]
  Seja $(X_n)_{n \geq 0}$ uma cadeia de Markov com espaço de estados
  $S$, distribuição inicial $\lambda$ e matriz de transição $P$.  Seja
  $T$ um tempo de parada para $(X_n)_{n \geq 0}$.  Então, condicionado
  em $T < \infty$ e $X_T = i$, o processo \( (X_{T+n})_{n \geq 0}\) é
  uma cadeia de Markov com estado inicial $i$ e matriz de transição
  $P$.  Além disso, $(X_{T+n})_{n \geq 0}$ é independente de
  $(X_0, X_1, \dots, X_T)$.
\end{theorem}



Sejam $A$ um evento que depende apenas do futuro
$(X_{T+1},X_{T+2},\dots)$ e $B$ um evento que depende apenas do passado
$(X_0,\dots,X_T)$. Pela definição de probabilidade condicional, temos
\[
\frac{\mathbb{P}(A \cap B \mid T<\infty, X_T=i)}{\mathbb{P}(B \mid T<\infty, X_T=i)}
= \mathbb{P}(A \mid B, T<\infty, X_T=i).
\]
A Propriedade (``fraca'') de Markov em tempo fixo garante que, dado
$X_T=i$, o futuro é independente do passado. Assim,
\[
\mathbb{P}(A \mid B, T<\infty, X_T=i) 
= \mathbb{P}(A \mid T<\infty, X_T=i).
\]
Por fim, a Propriedade Forte de Markov afirma que, condicionado em
$X_T=i$, o futuro $(X_{T+1},X_{T+2},\dots)$ tem a mesma lei que uma
cadeia de Markov iniciada em $i$. Ou seja,
\[
\mathbb{P}(A \mid T<\infty, X_T=i) = \mathbb{P} (A \mid X_0=i).
\]
Isso mostra que o futuro é independente do passado, e que o processo
reinicia como uma nova cadeia de Markov em $i$.

O Teorema~1 diz que, se \(T < \infty\), então, dado que \(X_T = i\):
\begin{enumerate}[label=(\roman*), noitemsep]
\item \((X_{T+1}, \ldots, X_{T+n})\) é independente de
  \((X_0, \ldots, X_{T-1})\); e
\item tem a mesma distribuição que \((X_1, \ldots, X_n)\) sob
  $\mathbb P(\cdot\mid X_0=i)$.
\end{enumerate}


\begin{proof}
  Sejam $i,j_1,\dots,j_n \in S$ e seja $B$ um evento determinado por
  $(X_0,\dots,X_T)$.  Nosso objetivo é mostrar que
  \begin{align*}
    \mathbb{P} ( & X_{T+1}=j_1,\dots,X_{T+n}=j_n,\, B \,\mid \,
                   T<\infty, X_T=i )\\ & = \mathbb{P}  (   X_{1}=j_1,\dots,X_n=j_n \mid  X_0=i)\cdot 
                                         \mathbb{P}(B \mid T<\infty, X_T=i).
  \end{align*}


\noindent\textbf{1.~Redução ao caso $T=m$.}  
Se $T=m$, o evento $B \cap [T=m]$ depende apenas de $(X_0,\dots,X_m)$.
No evento $[T=m]$ temos $X_T = X_m$, logo para cada $m$
\[
  \big[ X_T=i, X_{T+1}=j_1,\dots,X_{T+n}=j_n,\, B,\, T=m \big] \;=\;
  [ X_m=i ] \cap A_m  \cap B \cap [T=m]
\]
onde $A_m=[X_{m+1}=j_1,\dots,X_{m+n}=j_n]$.

Como $B \cap [T=m]$ depende apenas do passado até $m$, a propriedade de
Markov em $m$ implica em
\( \mathbb{P}(A_m \mid X_m=i, B, T=m) = \mathbb{P}(A_m \mid X_m=i) \),
ou seja, a probabilidade do futuro partindo do estado $i$, logo pela
definição de probabilidade condicional
\begin{align*}
\mathbb{P}(X_m=i, A_m, B, T=m) 
&\;=\; \mathbb{P}(B, T=m, X_m=i)\cdot \mathbb{P}(A_m \mid X_m=i, B, T=m). \\
&\;=\; \mathbb{P}(B, T=m, X_m=i)\cdot  \mathbb{P}(A_m \mid X_m=i).
\end{align*}
Portanto,
\begin{align}\notag
\mathbb{P}(X_T=i,\dots,X_{T+n}=j_n, B, T=m)
\;=\;\mathbb{P}(X_m=i, A_m, B, T=m) 
\\=
\label{eq:1}
  \mathbb{P}(B, T=m, X_m=i)\cdot
  \underbrace{\mathbb{P}(X_m=i,\dots,X_{m+n}=j_n \mid X_m=i)}_{=\mathbb{P}(X_0=i,\dots,X_{n}=j_n \mid X_0=i)}.
\end{align}

\medskip

\noindent\textbf{2.~Soma sobre os valores possíveis de $T$.}  
Os eventos $[T=m]$, para $m=0,1,2,\dots$, são disjuntos e cobrem
$[T<\infty]$.  Logo,
\begin{align*}
&\mathbb{P}(X_T=i,\dots,X_{T+n}=j_n, B, T<\infty) \\
&= \sum_{m=0}^\infty \mathbb{P}(X_T=i,\dots,X_{T+n}=j_n, B, T=m).
\end{align*}
Substituindo  no resultado acima
\begin{align*}
\mathbb{P}&(X_T=i,\dots,X_{T+n}=j_n, B, T<\infty) \\
&=\mathbb{P}(X_0=i,\dots,X_n=j_n \mid X_0=i)\cdot \sum_{m=0}^\infty \mathbb{P}(B, T=m, X_m=i) \\
&=\mathbb{P}(X_0=i,\dots,X_n=j_n \mid X_0=i)\cdot\mathbb{P}(B, T<\infty, X_T=i).
\end{align*}

\medskip

\noindent\textbf{3.~Condicionamento final.}  
Dividindo por $\mathbb{P}(T<\infty, X_T=i)$ concluímos:
\begin{align*}
\mathbb{P} & (X_T=i,\dots,X_{T+n}=j_n, B \mid T<\infty, X_T=i) 
\\  & = \mathbb{P}(X_0=i,\dots,X_n=j_n\mid X_0=i)\cdot \mathbb{P}(B \mid T<\infty, X_T=i).
\end{align*}


Isto prova que, condicionado em $T < \infty$ e $X_T = i$, o futuro
$(X_{T+n})_{n \geq 0}$ tem a mesma lei que uma cadeia de Markov
iniciada em $i$ e é independente do passado $(X_0,\dots,X_T)$.
\end{proof}






\end{document}
